\section{State Schema and Serialization Listings}
\label{app:state-listings}

This appendix collects the implementation-level listings for the
structured context state of Section~\ref{sec:method-state} and the
serialization template of Section~\ref{sec:method-architecture}. The
load-bearing definitions (Definition~\ref{def:mathstate}, the
fact-versus-hypothesis distinction) remain in the body; what follows is
the concrete record layout and token template.

\paragraph{Record dataclasses.}
Each record type in Definition~\ref{def:mathstate} carries a stable
identifier so that actions and parent pointers can address it. The
forward-compatible \texttt{parent\_*} and \texttt{children} fields are
left at their defaults in phase 1 and activated in phase 2 for
structural decomposition (Section~\ref{sec:method-state}).

\begin{lstlisting}
@dataclass(frozen=True)
class Fact:
    id: str
    statement: str
    source: Literal["given", "derived", "verified"]
    derivation: str | None

@dataclass(frozen=True)
class Subgoal:
    id: str
    statement: str
    parent_subgoal: str | None = None
    children: tuple[str, ...] = ()

@dataclass(frozen=True)
class Hypothesis:
    id: str
    statement: str
    parent_subgoal: str | None
    parent_hypothesis: str | None = None
    children: tuple[str, ...] = ()

@dataclass(frozen=True)
class VerifierCall:
    id: str
    verb: Literal["check_arith", "solve", "simplify",
                  "factor", "evaluate"]
    input: str
    output: str
    success: bool
    follow_up_to: str | None = None
\end{lstlisting}

\paragraph{Serialization template.}
A math state $s$ is rendered into a token sequence by the deterministic
template below, using the explicit structural markers discussed in
Section~\ref{sec:method-architecture}. Two properties of this template
matter and are explained there: the markers give the LLM's
self-attention stable positional anchors, and depth-indexed identifiers
render the implicit tree structure as a flat token sequence that causal
self-attention can resolve without a tree-aware encoder.

\begin{lstlisting}[basicstyle=\ttfamily\footnotesize]
[PROBLEM] {problem text} [/PROBLEM]
[FACT id=fact_1 source=given]      {statement} [/FACT]
[FACT id=fact_2 source=verified]   {statement} [/FACT]
[OPEN_SUBGOAL id=sg_1]             {statement} [/OPEN_SUBGOAL]
[HYPOTHESIS id=hyp_1 parent=sg_1  depth=1] {statement} [/HYPOTHESIS]
[HYPOTHESIS id=hyp_2 parent=hyp_1 depth=2] {statement} [/HYPOTHESIS]
[VERIFIER_CALL id=vcall_1 verb=factor success=true]
   in:  {input}
   out: {output}
[/VERIFIER_CALL]
[COMMITTED_ANSWER] <none> [/COMMITTED_ANSWER]
[DEPTH] {d} [/DEPTH]
[STATE_END]
\end{lstlisting}

\section{Action-Space Alternatives and Migration Path}
\label{app:action-alternatives}

The body (Section~\ref{sec:method-actions}) adopts the categorical
seven-verb action set (Definition~\ref{def:actionset}) as the
conservative phase 1 design. Two alternative action-space designs recur
in the LLM-MCTS literature and represent different bets.

\begin{description}[leftmargin=*, style=nextline]
  \item[Verb $\times$ slot.]
  Actions are pairs $(v, s)$ where $v$ is a verb and $s$ is a target
  slot in the state. The action space grows as
  $|V| \times |\text{slots}|$. This is more expressive for very
  structured domains, but with our small verb set and id-based
  targeting, the additional cardinality is unnecessary: the
  verb-instantiator already supplies the target identifier as part of
  its content, achieving the same expressivity without inflating the
  categorical action space.
  \item[Progressive widening with LLM-proposed candidates.]
  At each decision node, sample $K$ candidate actions from the LLM as
  free-form text; the children grow as
  $K(s) = \lceil \alpha \cdot N(s)^{\beta} \rceil$. This is the
  pattern in AlphaProof \citep{deepmind2024alphaproof} and AlphaCode
  \citep{li2022alphacode}. It is powerful when the action space is
  genuinely open (Lean tactics with arguments; arbitrary code
  completions) but requires an $(s, a)$-embedding-based prior in place
  of a categorical $Q$-head, and double progressive widening machinery
  to handle the growing branching factor.
\end{description}

Migration to progressive widening is planned for phase 2 once the
categorical version is stable; the architectural change required is
bounded (replace the categorical $Q$-head with an
$(s, a)$-embedding scorer), and the rest of the system
(serialization, adapter, verifier-anchored reward) carries over
unchanged.
